% WAIT -- os/samir/cmd/proc.c p_wait (prompt, one key via conin_char, a
% one-character string; Enter alone -> "").  Gate: test-interp-proc.
\begin{XbaseEntry}{WAIT}{command}
Waits for one key to be pressed, and can store it in a memory variable.
\Syntax
\SyntaxLine{\kw{WAIT} \opt{\xp{expC}} \opt{\kw{TO} \xp{memvar}}}
\Usage
\cmd{WAIT} displays the prompt \xp{expC}, if one is given, and waits until a key
is pressed.  The key is not followed by \key{Enter}.  With \kw{TO}, the key is
stored in \xp{memvar} as a character string of one character; if the key was
\key{Enter}, the memory variable is an empty string.  Without a prompt,
\cmd{WAIT} displays nothing while it waits.
\Tips
A program uses \cmd{WAIT} to hold its output on the screen until the operator
has read it, and to take a single-letter answer to a question.
\Examples
\DotExample{examples/c-wait}
\SeeAlso{ACCEPT, INPUT}
\end{XbaseEntry}

